\section{Exact topology, coercivity, and Gaussian separation}\label{sec:linear}

Let $D$ be the $216\times136$ integer matrix obtained by restricting $d_1$ to
the cotree links, and set
\begin{equation}\label{eq:H-def}
 H=D^{\mathsf T}D.
\end{equation}

\begin{proposition}[Exact dimensions]\label{prop:dimensions}
The cellular ranks and retained rank are
\[
 \rank d_0=80,\qquad \rank d_1=136,\qquad \rank R=17.
\]
Consequently
\[
 \dim\mathcal Y_0=17,\quad \dim F_0=119,
 \qquad \dim\cY=51,\quad \dim\cF=357.
\]
\end{proposition}

\begin{proof}
The one-skeleton is connected, so $\rank d_0=81-1=80$.  Contractibility of
$I_2^4$ gives $\ker d_1=\operatorname{im}d_0$, hence $\rank d_1=216-80=136$.
The face-sum map factors through the curvature map of one coarse four-cube,
whose rank is $32-(16-1)=17$.  Exact rational elimination on the displayed
rows in \eqref{eq:selected-rows} gives the matching lower bound after tree
gauge.  Rank--nullity and tensoring by $\R^3$ finish the proof.  The explicit
integer matrices appear in Appendix~\ref{app:linear-algebra}.
\end{proof}

\begin{theorem}[Certified cotree coercivity]\label{thm:cotree-coercivity}
For every $a\in\R^{136}$,
\begin{equation}\label{eq:cotree-floor}
 a^{\mathsf T}Ha=\norm{Da}^2\ge\frac1{25}\norm a^2.
\end{equation}
The same bound holds on $F_0$ and after tensoring with $\R^3$.
\end{theorem}

\begin{proof}
The supplement constructs $H$ from the declared cubical orientation and proves
that $H-\frac1{25}\Id$ admits an exact rational $LDL^{\mathsf T}$
factorization with $136$ positive pivots.  This is Sylvester positivity and
implies \eqref{eq:cotree-floor}.  A second implementation proves positivity of
$25H-\Id$ by fraction-free Bareiss leading-principal-minor elimination.  No
floating eigenvalue is used.  See Appendix~\ref{app:linear-algebra}.
\end{proof}

We now choose coordinates adapted to conditioning.  Put
\[
 G=RH^{-1}R^{\mathsf T},\qquad
 \mathsf B=G^{-1/2}R,\qquad M=H^{-1}\mathsf B^{\mathsf T}.
\]
Choose an Euclidean isometry $N:\R^{119}\to\ker R$.  Then
\begin{equation}\label{eq:horizontal-identities}
 \mathsf BM=\Id_{17},\quad M^{\mathsf T}HM=\Id_{17},\quad
 M^{\mathsf T}HN=0,\quad N^{\mathsf T}N=\Id_{119},
 \quad \norm M\le5.
\end{equation}
The last inequality follows from Theorem~\ref{thm:cotree-coercivity}.  We use
the same letters for the color-tensored maps.  Every reduced logarithmic field
has the unique form
\begin{equation}\label{eq:decomposition}
 A=My+Nz,\qquad y\in\R^{51},\quad z\in\R^{357}.
\end{equation}

\begin{proposition}[Exact centered Gaussian cancellation]\label{prop:gaussian}
For the tangent density proportional to
$\exp[-A^{\mathsf T}(H\otimes\Id_3)A/(2t)]$, conditioning on $\mathsf BA=y$ gives
\[
 A=My+NZ,\qquad
 Z\sim\mathcal N\!\left(0,t(N^{\mathsf T}HN)^{-1}\right).
\]
In particular, the centered fiber law is exactly independent of $y$.
\end{proposition}

\begin{proof}
Insert \eqref{eq:decomposition} into the quadratic form and use
\eqref{eq:horizontal-identities}.  The cross term vanishes and the retained
quadratic form is $\norm y^2$.  Normalization produces the stated Gaussian.
\end{proof}
